\documentclass[10pt,a4paper]{amsart}
\usepackage[margin=19mm]{geometry}
\usepackage{amsmath,amssymb,mathtools}
\usepackage[scr=rsfs,cal=euler]{mathalfa}
\usepackage{xfrac}
\usepackage[unicode,hidelinks,bookmarks=false]{hyperref}
\newcommand{\ie}{i.e.\ }
\newcommand{\cf}{cf.\ }
\newcommand{\resp}{respectively\ }
\newcommand{\Subsection}{\subsection}
\providecommand{\nobreakdash}{\nobreak\hbox{-}}
\setlength{\emergencystretch}{2em}
\multlinegap0pt
\theoremstyle{plain}
\newtheorem{lemma}[equation]{Lemma}
\newtheorem{theorem}[equation]{Theorem}
\newtheorem{proposition}[equation]{Proposition}
\newtheorem{corollary}[equation]{Corollary}
\newtheorem{claim}[equation]{Claim}
\newtheorem{defn}[equation]{Definition}
\theoremstyle{remark}
\newtheorem{remark}[equation]{Remark}
\title[Local points on twists of X(p)]{Local points on twists of $X(p)$ with applications: two-dimensional mod p representations with equal determinants and equal inertia restrictions}
\author{Nuno Freitas, Diana Mocanu}
\date{Source: arXiv:2509.04294v4 (CC BY 4.0)}
\newcommand{\F}{\mathbb{F}}
\DeclareMathOperator{\GL}{GL}
\newcommand{\Gal}{\text{Gal}}
\newcommand{\Q}{\mathbb{Q}}
\newcommand{\Z}{\mathbb{Z}}
\newcommand{\rhobar}{{\overline{\rho}}}
\begin{document}
\maketitle
\noindent\emph{Source and licence.} Local points on twists of $X(p)$ with applications. Version arXiv:2509.04294v4 (CC BY 4.0), distributed under the Creative Commons Attribution 4.0 International licence, \url{http://creativecommons.org/licenses/by/4.0/}, as recorded in the arXiv licence metadata for this exact version and shown on the abstract page \url{https://arxiv.org/abs/2509.04294v4}. Full licence notice: licenses/arxiv-2509.04294-CC-BY-4.0.txt.\par\smallskip

\noindent\textit{Editorial scope.} Counted unit: the article's Lemma with source label \texttt{lem:niveau-two-representations} (source0.tex lines 2411--2423), the rigidity statement that two continuous two-dimensional mod p representations of the absolute Galois group of the p-adics with equal determinants and isomorphic restrictions to inertia, of a non-degenerate niveau two shape, are isomorphic; delivered together with the article's complete original proof of it (source0.tex lines 2425--2484, one \texttt{proof} environment of 60 lines, adjacent to the statement). No mathematics is written, repaired or re-derived: statement and proof are the article's own text line for line, in the article's own notation (the Galois groups, the inertia subgroup, the niveau two characters, the scalar extensions and the matrices), and no retained line is substituted, re-flowed or omitted. No further part of the article is required as context and none is retained: the counted proof uses only the objects introduced in its own statement, so no cross-reference and no citation occurs in any retained line; consequently the wrapper carries no bibliography and no bibliography entry, and no context passage is delivered. Nothing else of the article is retained: its main theorems, its other lemmas and its applications are omitted. Editorial formatting only: an AMS \texttt{amsart} preamble that restates the environments the retained lines use, namely the article's declarations of Lemma, Theorem, Proposition, Corollary, Claim, Definition and Remark, all sharing the article's equation counter as in the article, and the macros that the retained lines need. The article's own source declares the class \texttt{amsart} as well; no class or style file is shipped with this package. Numbering differs from the article, because the article ties its theorem-like counters to the equation counter and because only one environment is retained: the counted lemma prints with the wrapper's own counter. No picture environment and no graphical inclusion occurs in any retained line, and the package delivers no image asset. One bundled source file, \texttt{sources/184479/source0.tex}, accompanies the delivered item as the exact source of every retained span. Every retained span is recorded in \texttt{delivery-evidence.json} with method \texttt{exact\_tex}.
\begin{lemma}\label{lem:niveau-two-representations}
Let $p > 2$ be a prime. Let
$\rho_1,\rho_2:G_{\Q_p}\longrightarrow \GL_2(\F_p)$
be continuous representations with equal determinants.
Suppose that, for some  $k \in \Z$, we have $\omega_2^k\neq\omega_2^{pk}$ and
\[
 \left.\rho_i\right|_{I_p}\otimes_{\F_p}\F_{p^2}
 \simeq
 \omega_2^k\oplus\omega_2^{pk},
 \qquad i=1,2.
\]
Then $\rho_1\simeq\rho_2$.
\end{lemma}

\begin{proof}
Set $\theta=\omega_2^k$. By assumption, we have $\theta\neq\theta^p$ and, after extending scalars to $\F_{p^2}$, there is a basis such that, for all $\sigma \in I_p$, we have
\[
 \rho_i(\sigma)=
 \begin{pmatrix}
  \theta(\sigma)&0\\
  0&\theta^p(\sigma)
 \end{pmatrix}.
\]
Let $\sigma\in I_p$ be a generator of the tame inertia
quotient, and let $\tau\in G_{\Q_p}$ be a lift of arithmetic
Frobenius. Recall the tame relation
$\tau\sigma\tau^{-1}=\sigma^p$. Since $\theta \neq \theta^p$, it follows that conjugation by $\rho_i(\tau)$ exchanges the 
eigenspaces corresponding to $\theta$ and $\theta^p$. Therefore,
\[
 \rho_i(\tau)=
 \begin{pmatrix}
  0&a_i\\
  b_i&0
 \end{pmatrix}
\]
for some $a_i,b_i\in\F_{p^2}^{\times}$. From the equality of determinants, we have $a_1b_1=a_2b_2$.
It follows that the matrix
$
 D=\left(
 \begin{smallmatrix}
  a_2/a_1&0\\
  0&1
 \end{smallmatrix}\right)
$ 
centralizes image of inertia and satisfies
$D\rho_1(\tau)D^{-1}=\rho_2(\tau)$.
Thus
\[
 D\rho_1(g) =  \rho_2(g) D, \quad \text{ for all } g \in G_{\Q_p},
\]
showing the representations are isomorphic over $\F_{p^2}$.
It remains to descend this isomorphism to $\F_p$. 
Let $P_i\in\GL_2(\F_{p^2})$ be the matrix that transforms $\rho_i$ from the original
$\F_p$-basis to the matrix description above. Then
 $A:=P_2DP_1^{-1}\in\GL_2(\F_{p^2})$ satisfies
\[
 A\rho_1(g)=\rho_2(g)A
 \qquad\text{for all }g\in G_{\Q_p}
\]
where now $\rho_i(g) \in \GL_2(\F_p)$ for all $g \in G_{\Q_p}$.
Let $\delta$ be the nontrivial element of
$\Gal(\F_{p^2}/\F_p)$. We also have  $\delta(A)\rho_1(g)=\rho_2(g)\delta(A)$ 
for all $g\in G_{\Q_p}$. 
Note that the matrix description above also shows the $\rho_i$ are absolutely irreducible. Therefore, by Schur's
lemma, there exists $c\in\F_{p^2}^{\times}$ such that
$\delta(A)=cA$. Applying $\delta$ again gives
$c\,\delta(c)=1$. By Hilbert's Theorem~90, there is
$\lambda\in\F_{p^2}^{\times}$ such that
$c=\frac{\lambda}{\delta(\lambda)}$.
Consequently,
$\delta(\lambda A)=\lambda A$,
so that $\lambda A\in\GL_2(\F_p)$. Hence $\lambda A$ gives an
isomorphism $\rho_1\simeq\rho_2$ over $\F_p$.
\end{proof}

\end{document}
